\chapter{Pre-Study Administrative Process Validation Survey}
\label{app:pre-admin-survey}

This survey is for university personnel involved in processing applications for AdDU's internal scholarship tracks. Please answer based on the most recent application cycle handled through the manual, paper-based process. Please complete this survey before seeing or testing ScholarPath AdDU.

Question 1: What is your role in the scholarship application process?

\begin{itemize}
  \item Question Type: Multiple Choice
  \item Options:
  \begin{itemize}
    \item Office head or supervisor
    \item Staff (screening and document verification)
    \item School Scholarship Subcommittee member
    \item Other (please specify) Component 1: Perceived Processing Time Question 2: On average, how much staff time does it take to process a single scholarship application, from receiving it to marking it ready for interview or subcommittee evaluation? (Include screening and document verification; exclude the interview itself.)
  \end{itemize}
  \item Question Type: Multiple Choice
  \item Options:
  \begin{itemize}
    \item Less than 15 minutes
    \item 15 -- 30 minutes
    \item 31 -- 60 minutes
    \item 1 -- 2 hours
    \item More than 2 hours Question 3: How long does one full application cycle take, from the opening of applications to the release of results?
  \end{itemize}
  \item Question Type: Multiple Choice
  \item Options:
  \begin{itemize}
    \item Less than 2 weeks
    \item 2 -- 4 weeks
    \item 1 -- 2 months
    \item 2 -- 3 months
    \item More than 3 months Component 2: Perceived Effort Scale Question 4: On a scale of 1 to 5, please rate how easy or difficult each of the following tasks is for your office under the manual process.
  \end{itemize}
  \item Question Type: Likert Grid (5-point scale)
  \item Scale Columns:
  \begin{itemize}
    \item 1 - Very Easy
    \item 2 - Easy
    \item 3 - Neutral
    \item 4 - Difficult
    \item 5 - Very Difficult
  \end{itemize}
  \item Row Sub-tasks: 1. Receiving, sorting, and filing incoming paper applications. 2. Checking whether each applicant meets the baseline requirements (high school average, entrance exam result, household income) of the chosen track. 3. Verifying submitted documents against official report cards and income documents. 4. Requesting missing or corrected documents from applicants and tracking whether they were submitted. 5. Collating evaluation results from interview panels and the five School Scholarship Subcommittees. 6. Identifying who made or approved a specific status change or decision on an application. 7. Informing applicants of their application status, correction requests, and deadlines. Component 3: Processing Issues History Question 5: In a typical application cycle, how often does each of the following happen?
  \item Question Type: Likert Grid (5-point frequency scale)
  \item Scale Columns:
  \begin{itemize}
    \item 1 - Never
    \item 2 - Rarely
    \item 3 - Sometimes
    \item 4 - Often
    \item 5 - Very Often
  \end{itemize}
  \item Row Items: 1. A full application file is processed even though the applicant does not meet the baseline requirements. 2. An application is incomplete or contains an illegible document. 3. An applicant does not submit requested missing or corrected documents before the deadline. 4. An applicant misses a submission deadline. 5. Interview or subcommittee results have to be re-encoded or collated manually. 6. It is unclear who made or approved a decision on an application. 7. Applicants contact the office to ask about the status of their application. Question 6 (Optional): Which part of the manual process creates the most work for your office? Please describe briefly.
  \item Question Type: Short Answer (Open-ended)
\end{itemize}
